\def\topic{Measure theory}
\ifdefined\MyHeaderIncluded
\else
\def\qwertyflag{true}

\ifdefined\MyHeaderIncluded

\else
  \newcommand{\foobar}{FooBar}
  \def\MyHeaderIncluded{true}
% if is till the end of the file

\documentclass[10pt,a4paper]{article}
\usepackage[utf8]{inputenc}
%\usepackage[german]{babel}
%\usepackage[cyrillic]{babel}
\usepackage[T1]{fontenc}
%\usepackage[T2A]{fontenc}
%\usepackage[english, russian]{babel}
\usepackage{amsmath}
\usepackage{amsfonts}
\usepackage{amssymb}
\usepackage{amsthm}
\usepackage{mathtools} %coloneqq

\usepackage{bussproofs} %

%\renewcommand{\thesubsection}{\arabic{subsection}}
%\newtheorem{thm}{Theorem}[section]
%\renewcommand{\thethm}{\arabic{subsection}.\arabic{thm}}

\usepackage{thmtools}
\usepackage{hyperref}
\hypersetup{
    colorlinks=true,
    linkcolor=blue,
    filecolor=magenta,      
    urlcolor=cyan,
    pdftitle={Overleaf Example},
    %pdfpagemode=FullScreen,
    }
%\usepackage{stix-type1}%\turnediota
\theoremstyle{definition}

\newtheorem{constructor}{Constructor}[section]
\newtheorem{iaxiom}{Axiom}
\newtheorem{axiom}{Axiom}
\newtheorem*{axiom*}{Axiom}
\newtheorem{definition}{Definition}[section]
\newtheorem*{definition*}{Definition}
\newtheorem{remark}{Remark}[section]
\newtheorem*{remark*}{Remark}
\newtheorem{recall}{Recall}[section]
\newtheorem*{recall*}{Recall}

\newtheorem{lemma}{Lemma}[section]
\newtheorem{theorem}{Theorem}[section]
\newtheorem*{theorem*}{Theorem}
\newtheorem{note}{Note}[section]
\usepackage{etoolbox}

%\newcommand*{\@Dashv}[2]{%
%  \reflectbox{$\m@th#1#2$}%
%}
\newcommand*{\separation}[3]{% 1)binder, 2)set, 3)condition
  \{#1\in #2: #3\}
}

% https://www.freecodecamp.org/news/simple-regex-tricks-for-beginners-3acb3fa257cb/
%  \\mathcal\{T\} |---> \topology
%  \\mathcal\{([^\}]*)\}   |--->  $1
\newcommand{\AND}{\ \ \&\ \ }
\newcommand{\any}{\mathrm{\ any}}
\newcommand{\borel}{\mathcal{B}}
\newcommand{\can}{\mathrm{can}}
\newcommand{\canonize}{\mathrm{canonize}}
\newcommand{\closure}{\mbox{cl}}
\newcommand{\complex}{\mathbb{C}}
\newcommand{\converges}{\longrightarrow}
\newcommand{\countable}{\mbox{\ countable}}
\newcommand{\ldefined}{\mbox{\ defined}}
\newcommand{\dom}{\mbox{dom}}
\newcommand{\Dom}{\mbox{Dom}}
\newcommand{\Done}{\noindent \textbf{Done.}\\}
\newcommand{\lexists}{\mbox{\ exists}}
\newcommand{\lforall}{\mbox{\ forall}}
\newcommand{\lunique}{\mbox{\ unique}}
\newcommand{\family}{\mathrm{\ family}}
\newcommand{\Fm}{\mathrm{Fm}}
\newcommand{\FuncSymb}{\mathrm{FuncSymb}}
\newcommand{\function}{\mathrm{\ function}}
\newcommand{\functional}{\mbox{\ functional}}
\newcommand{\IFF}{\ \Longleftrightarrow\ }
\renewcommand{\iff}{\ \longleftrightarrow\ }
\newcommand{\IMPLIES}{\ \Longrightarrow\ }
\renewcommand{\implies}{\ \longrightarrow\ }
\newcommand{\Inf}{\mathrm{Inf}}
\newcommand{\injective}{\mbox{\ injective}}
\newcommand{\injection}{\mbox{\ injection}}
\newcommand{\interior}{\mbox{int}}
\newcommand{\intersection}{\bigcap}
\newcommand{\isempty}{\mathrm{\ empty}}
\newcommand{\LB}{\mathrm{LB}}
\newcommand{\nghbset}{\mathcal{N}}
\newcommand{\nonempty}{\mathrm{\ nonempty}}
\newcommand{\Ord}{\mathrm{Ord}}
\newcommand{\power}{\mathcal{P}}
\newcommand{\pr}{\mathrm{pr}}
\newcommand{\fst}{\mathrm{fst}}
\newcommand{\snd}{\mathrm{snd}}
\newcommand{\PredSymb}{\mathrm{PredSymb}}
\newcommand{\Proof}{\noindent \textbf{Proof.}\\}
\newcommand{\ran}{\mathrm{ran}}
\newcommand{\Ran}{\mathrm{Ran}}
\newcommand{\real}{\mathbb{R}}
\newcommand{\relation}{\mathrm{\ relation}}
\newcommand{\set}{\mathrm{\ set}}
\newcommand{\singleton}{\mathrm{\ singleton}}
\newcommand{\semantics}{\mathrm{semantics}}
\newcommand{\semialgebra}{\mathcal{S}}
\newcommand{\Sup}{\mathrm{Sup}}
\newcommand{\surjective}{\mbox{\ surjective\ on\ }}
\newcommand{\Term}{\mathrm{Term}}
\newcommand{\Type}{\mathrm{Type}}
\newcommand{\topology}{\mathcal{T}}
\newcommand{\total}{\mbox{\ total\ on\ }}
\newcommand{\transitive}{\mathrm{\ transitive}}
\newcommand{\union}{\bigcup}
\newcommand{\lTrue}{\mathrm{True}}
\newcommand{\lFalse}{\mathrm{False}}
\newcommand{\lUncertain}{\mathrm{Uncertain}}
\newcommand{\lif}{\mathrm{if}}
\newcommand{\lthen}{\mathrm{then}}
\newcommand{\lelse}{\mathrm{else}}
\newcommand{\lexception}{\mathrm{exception}}
\newcommand{\lundef}{\mathrm{undef}}
\newcommand{\univocal}{\mathrm{\ univocal}}
\newcommand{\univocalat}{\mathrm{\ univocal\ at\ }}
\newcommand{\UB}{\mathrm{UB}}
\newcommand{\Un}{\mathcal{U}}
\newcommand{\Var}{\mathrm{Var}}

\newcommand{\notvarphi}{\phi}
\renewcommand{\phi}{\varphi}

\newcommand{\comment}[1]{\quad\mbox{(#1)}}
% local commands
%\preto\section{\undef\blipo} %etoolbox
\preto\subsection{\undef\blipo} %etoolbox
% https://tex.stackexchange.com/questions/105684/limit-macro-newcommand-scope-to-section
\usepackage{graphicx}
\usepackage{authblk}

\setlength\parindent{0pt} % whole doc of \noindent
\usepackage{soul}
%\usepackage{mathgl}

\fi % guard
\ifdefined\topic
\else
\def\topic{No Titel}
\fi
\title{\topic}
\author[1]{Georgy Dunaev}
\begin{document}
\maketitle
\fi
\section{Measure theory}
The functional-analytic applications require a specified measure space,
not merely a collection called a sigma-algebra. The following conventions
connect the earlier closure axioms with spaces of integrable functions.
The Lebesgue integral is constructed below from nonnegative simple
functions. Its convergence theorems then support the completeness,
duality, and compactness arguments here and in the following chapter.

\subsection{Carrier and measure}
Fix a set $\Omega$ and a sigma-algebra $\mathcal{A}\subseteq\power\Omega$
with $\Omega\in\mathcal{A}$. Thus the carrier is explicit; it also equals
$\union\mathcal{A}$, in agreement with the earlier convention.
A positive measure is a function
\[
\mu:\mathcal{A}\longrightarrow[0,+\infty]
\]
with $\mu(\emptyclass)=0$ and
\[
\mu\left(\bigcup_{n=1}^{\infty}A_n\right)
 =\sum_{n=1}^{\infty}\mu(A_n)
\]
for every sequence of pairwise disjoint members of $\mathcal{A}$.
The measure is finite when $\mu(\Omega)<\infty$.
Finiteness is an additional hypothesis, not a consequence of the carrier
being a set.

A statement holds \emph{almost everywhere} if it fails only on a subset
of a measurable set of measure zero. When the measure is not complete,
that subset need not itself belong to $\mathcal{A}$. Measurability
requirements must therefore still be stated separately.

\subsection{Generating classes and uniqueness of finite measures}
All families in this subsection are families of subsets of the same
fixed set $\Omega$. A \emph{$\pi$-system} is a family closed under
binary intersections; we permit the empty family. A \emph{Dynkin
system}, or \emph{$\lambda$-system}, is a family $\mathcal D$ that
contains $\Omega$, contains $\Omega\setminus A$ whenever it contains
$A$, and contains the union of every countable pairwise disjoint
sequence of its members. In particular $\emptyclass\in\mathcal D$.
For nested $A\subseteq B$ in $\mathcal D$, it also contains
\[
B\setminus A
=\Omega\setminus\bigl(A\cup(\Omega\setminus B)\bigr),
\]
since the two sets inside the union are disjoint. This recovers the
nested-difference formulation used in the topology chapter; conversely,
that formulation gives complements by taking $B=\Omega$.

Write $\sigma(\mathcal P)$ for the smallest sigma-algebra on
$\Omega$ containing a family $\mathcal P$, and $\lambda(\mathcal P)$
for the smallest Dynkin system on $\Omega$ containing it. These are
intersections of all the corresponding families containing
$\mathcal P$. Such collections are nonempty, since $\power\Omega$
belongs to them. The earlier intersection results apply, with the
common carrier $\Omega$ present in every family.

\begin{theorem}[Dynkin's $\pi$--$\lambda$ theorem]
\label{thm:pi_lambda}
If $\mathcal P$ is a $\pi$-system of subsets of $\Omega$, then
$\lambda(\mathcal P)=\sigma(\mathcal P)$. Consequently, every
Dynkin system containing $\mathcal P$ contains $\sigma(\mathcal P)$.
\end{theorem}
\Proof
Put $\mathcal P_0=\mathcal P\cup\{\Omega\}$. This is still a
$\pi$-system, even when $\mathcal P$ is empty, and adjoining
$\Omega$ changes neither of the generated families. Let
$\mathcal D=\lambda(\mathcal P_0)$. For any fixed $A\in\mathcal D$
consider
\[
\mathcal D_A=\{B\in\mathcal D:A\cap B\in\mathcal D\}.
\]
This is a Dynkin system on $\Omega$. It contains $\Omega$ because
$A\cap\Omega=A$. If $B\in\mathcal D_A$, then
$\Omega\setminus B\in\mathcal D$ and
\[
A\cap(\Omega\setminus B)=A\setminus(A\cap B)\in\mathcal D
\]
by the nested-difference property. Thus it is closed under complements.
If $B_n\in\mathcal D_A$ are pairwise disjoint, their union is in
$\mathcal D$, and the sets $A\cap B_n$ are pairwise disjoint members
of $\mathcal D$. Their union is $A\cap\bigcup_nB_n$, which belongs
to $\mathcal D$. This proves the required disjoint-union closure.

First take $A\in\mathcal P_0$. Closure of $\mathcal P_0$ under
intersections gives $\mathcal P_0\subseteq\mathcal D_A$.
Minimality of $\mathcal D$ then gives
$\mathcal D\subseteq\mathcal D_A$. We have proved that
$A\cap B\in\mathcal D$ whenever $A\in\mathcal P_0$ and
$B\in\mathcal D$.

Next fix an arbitrary $B\in\mathcal D$. The preceding conclusion
and symmetry of intersection give
$\mathcal P_0\subseteq\mathcal D_B$. Since $\mathcal D_B$ is a
Dynkin system, minimality again gives
$\mathcal D\subseteq\mathcal D_B$. Thus $\mathcal D$ is closed
under intersections of any two of its members.

Complements and intersections now give finite unions and differences.
For any sequence $A_n\in\mathcal D$, the sets
\[
C_n=A_n\setminus\bigcup_{j<n}A_j
\]
belong to $\mathcal D$, are pairwise disjoint, and have union
$\bigcup_nA_n$. Disjoint-union closure therefore gives arbitrary
countable unions. Hence $\mathcal D$ is a sigma-algebra containing
$\mathcal P$, so $\sigma(\mathcal P)\subseteq\mathcal D$.
The reverse inclusion holds because every sigma-algebra is a Dynkin
system. Equality and the final assertion follow.
\Done

\begin{theorem}[Uniqueness from a generating $\pi$-system]
\label{thm:finite_measure_uniqueness}
Let $\mu,\nu$ be finite positive measures on
$(\Omega,\mathcal A)$ with $\mu(\Omega)=\nu(\Omega)$.
If $\mathcal P\subseteq\mathcal A$ is a $\pi$-system with
$\sigma(\mathcal P)=\mathcal A$ and
$\mu(P)=\nu(P)$ for every $P\in\mathcal P$, then $\mu=\nu$.
If $\Omega\in\mathcal P$, equality of the total masses is already
part of the agreement assumption.
\end{theorem}
\Proof
The agreement class
\[
\mathcal D=\{A\in\mathcal A:\mu(A)=\nu(A)\}
\]
contains $\Omega$ by the total-mass assumption. For $A\in\mathcal D$,
finiteness permits subtraction and gives
\[
\mu(\Omega\setminus A)=\mu(\Omega)-\mu(A)
=\nu(\Omega)-\nu(A)=\nu(\Omega\setminus A).
\]
For pairwise disjoint $A_n\in\mathcal D$, countable additivity gives
\[
\mu\left(\bigcup_nA_n\right)=\sum_n\mu(A_n)
=\sum_n\nu(A_n)=\nu\left(\bigcup_nA_n\right).
\]
Thus $\mathcal D$ is a Dynkin system containing $\mathcal P$.
The $\pi$--$\lambda$ theorem implies
$\mathcal A=\sigma(\mathcal P)\subseteq\mathcal D$, proving equality
on every measurable set. If $\mathcal P$ is empty, its generated
sigma-algebra is $\{\emptyclass,\Omega\}$ and the same argument
applies; when $\Omega$ is empty this family has just one member.
\Done

\begin{example}[An arbitrary generating family is insufficient]
Take $\Omega=\{1,2,3,4\}$ and
$A=\{1,2\}$, $B=\{1,3\}$. The family
$\mathcal G=\{\Omega,A,B\}$ generates $\power\Omega$ as a
sigma-algebra: its singleton atoms are
$A\cap B$, $A\setminus B$, $B\setminus A$, and
$\Omega\setminus(A\cup B)$. But $\mathcal G$ is not a
$\pi$-system, since $A\cap B=\{1\}\notin\mathcal G$.
Define probability measures on all subsets by
\[
\mu(C)=\tfrac12|C\cap\{1,4\}|,\qquad
\nu(C)=\tfrac12|C\cap\{2,3\}|.
\]
The bars denote finite cardinalities. These finite sums of point masses are
countably additive: a point belongs to at most one member of a
pairwise disjoint sequence. Both measures give $\Omega$ mass one
and give $A$ and $B$ mass $1/2$, but
$\mu(\{1\})=1/2$ and $\nu(\{1\})=0$.
Thus agreement on generators alone does not replace agreement on a
generating intersection-closed family.
\end{example}

The uniqueness theorem lets probability laws be compared on a smaller
class of testing sets, provided that class is a generating $\pi$-system.
The proof used finite total masses when taking complements; no
infinite-measure uniqueness assertion is being inferred from it.

\subsection{Measurable scalar functions}
The Borel sigma-algebra of a topological space is the sigma-algebra
generated by its open sets. A function from $(\Omega,\mathcal A)$ to
a Borel space is measurable when the inverse image of every Borel set
belongs to $\mathcal A$. It suffices to check a generating family:
the sets in the target whose inverse images belong to $\mathcal A$
form a sigma-algebra, since inverse images preserve complements and
countable unions.

For real-valued functions, and for functions with values in the extended
real line or its subspace $[0,\infty]$, measurability is equivalent to
\[
\{\omega:f(\omega)>q\}\in\mathcal A\qquad(q\in\mathbb Q).
\]
Here the extended line has its order topology. Rational rays generate
its Borel sigma-algebra: complements and countable unions give both
open rays and open intervals, including neighbourhoods of the infinite
endpoints; these intervals have a countable base. For complex-valued
functions, measurability is equivalent to measurability of the real
and imaginary parts, since rational open rectangles form a countable
base for $\mathbb R^2$.

These descriptions justify the closure operations used below. For a
sequence of extended-real measurable functions,
\[
\{\sup_n f_n>a\}=\bigcup_n\{f_n>a\},\qquad
\{\inf_n f_n<a\}=\bigcup_n\{f_n<a\}.
\]
Thus countable suprema and infima are measurable, as are lower and
upper limits by their iterated infimum and supremum formulas. In
particular, an everywhere existing real pointwise limit is measurable;
the same holds for complex limits by taking the two coordinates.
The pair of two finite scalar measurable functions is measurable:
inverse images of basic rectangles are intersections of the two
coordinate inverse images, and every open set is a countable union
of basic rectangles. Composition with a continuous scalar operation
therefore preserves measurability. This proves measurability of sums,
products, absolute values, real maxima and minima, and truncations
of finite scalar functions, without any boundedness assumption.

For nonnegative extended-real functions, sums and products are obtained
as increasing limits of the corresponding operations on finite
truncations; for products this uses $0\cdot\infty=0$. They too are
measurable. No operation $\infty-\infty$ is defined by these statements.
Finally, replacing a measurable function by zero on a measurable set
preserves measurability, by splitting each inverse image over that set
and its complement. In particular this applies to measurable null
exceptional sets; arbitrary modifications on a nonmeasurable subset
of a null set need not preserve measurability for an incomplete measure.

\subsection{Constructing the integral}
The construction in this subsection works for every positive measure
space; neither finite total measure nor completeness of the measure is
assumed. All functions mentioned as integrands are measurable.

\begin{lemma}[Continuity from below]
If $A_n\subseteq A_{n+1}$ are measurable, then
$\mu(\bigcup_nA_n)=\lim_n\mu(A_n)$, with infinity allowed.
\end{lemma}
\Proof
Write the union as the disjoint union of $A_1$ and
$A_{n+1}\setminus A_n$ for $n\geq1$. Countable additivity identifies its
measure with the sum of their measures; the $n$th partial sum is exactly
$\mu(A_n)$. No subtraction of infinite measures is involved.
\Done

\begin{lemma}[Continuity from above]
\label{lem:measure_continuity_above}
If measurable sets satisfy $A_{n+1}\subseteq A_n$ and
$\mu(A_1)<\infty$, then
\[
\mu\Bigl(\bigcap_{n\geq1}A_n\Bigr)=\lim_{n\to\infty}\mu(A_n).
\]
\end{lemma}
\Proof
The sets $B_n=A_1\setminus A_n$ increase to
$A_1\setminus\bigcap_n A_n$. Finite additivity and the finite measure
of $A_1$ give $\mu(B_n)=\mu(A_1)-\mu(A_n)$; the same identity holds
for their union and the complementary intersection. Apply continuity
from below and subtract the resulting finite quantities from
$\mu(A_1)$. This gives the stated limit without an expression of the
form $\infty-\infty$.
\Done

The finiteness hypothesis cannot be omitted: with counting measure on
$\mathbb N$, the decreasing tails $A_n=\{k:k\geq n\}$ all have
infinite measure, whereas their intersection is empty. On a finite
measure space the hypothesis is automatic, as used later when passing
from almost-everywhere convergence to convergence in measure.

For later use, countable subadditivity also follows from countable
additivity. Given measurable $A_n$, replace them by the disjoint sets
$B_n=A_n\setminus\bigcup_{k<n}A_k$, which have the same union.
Monotonicity of measure follows from disjoint decomposition of an inclusion;
thus $\mu(B_n)\leq\mu(A_n)$ and
$\mu(\bigcup_nA_n)=\sum_n\mu(B_n)\leq\sum_n\mu(A_n)$.

A nonnegative simple function is a measurable function with finitely many
finite nonnegative values. For a disjoint measurable partition and
$s=\sum_{j=1}^m a_j\mathbf1_{A_j}$, define
\[
\int s\,d\mu=\sum_{j=1}^m a_j\mu(A_j),\qquad 0\cdot\infty=0.
\]
A common refinement of two finite partitions proves independence of the
representation, additivity, nonnegative homogeneity, and monotonicity for
simple functions: on each cell these assertions are the corresponding
scalar identities or inequalities. For $f:\Omega\to[0,\infty]$, put
\[
\int f\,d\mu=\sup\{\int s\,d\mu:0\leq s\leq f,
\ s\text{ simple}\}.
\]
The zero function is admissible. Simple monotonicity shows that this
extends the simple integral; the definition immediately gives monotonicity
for arbitrary nonnegative integrands. Integrating over $A$ means integrating
$f\mathbf1_A$, where its value off $A$ is defined to be zero even if
$f$ takes the value infinity there.

\begin{theorem}[Monotone convergence]
\label{thm:monotone_convergence}
If $0\leq f_n\uparrow f$ pointwise, then
$\int f_n\,d\mu\uparrow\int f\,d\mu$.
\end{theorem}
\Proof
Let $L=\sup_n\int f_n$. Monotonicity gives $L\leq\int f$.
Fix a nonnegative simple $s\leq f$ and $0<c<1$. The measurable sets
$E_n=\{f_n\geq cs\}$ increase to $\Omega$: at a point where $s>0$,
the limiting value is at least $s>cs$, and where $s=0$ the inequality
holds for every $n$. Since $cs\mathbf1_{E_n}\leq f_n$, simple
monotonicity and the definition give
\[
L\geq c\int s\mathbf1_{E_n}\,d\mu.
\]
Apply continuity from below to the finitely many level sets of $s$
intersected with $E_n$. It follows that
$L\geq c\int s$. Let $c\uparrow1$ (the assertion is also valid if
$\int s=\infty$), and then take the supremum over $s$. Thus
$L\geq\int f$, proving equality.
\Done

Every nonnegative measurable $f$ admits an increasing simple approximation:
round $\min(f,n)$ down to the nearest multiple of $2^{-n}$, treating
$\min(\infty,n)=n$. These simple functions increase to $f$.
For $f,g\geq0$, choose these approximations $s_n,t_n$. Then
$s_n+t_n\uparrow f+g$, so simple additivity and monotone convergence give
\[
\int(f+g)\,d\mu=\int f\,d\mu+\int g\,d\mu.
\]
Likewise $\int cf=c\int f$ for $c\geq0$, with the same zero convention.
All these expressions avoid an undefined difference $\infty-\infty$.
Applying monotone convergence to partial sums also proves
$\int\sum_n f_n=\sum_n\int f_n$ for nonnegative measurable $f_n$.

\begin{lemma}[Null sets and finite integrals]
For $f\geq0$, $\int f=0$ if and only if $f=0$ almost everywhere.
If $\int f<\infty$, then $f<\infty$ almost everywhere. Nonnegative
measurable functions equal almost everywhere have equal integrals.
\end{lemma}
\Proof
For every $a>0$, comparison with $a\mathbf1_{\{f>a\}}$ gives
$a\mu(\{f>a\})\leq\int f$. If the integral is zero, the measurable
set $\{f>0\}=\bigcup_{n\geq1}\{f>1/n\}$ is null. Conversely, a
simple function below $f$ has every positive level supported on that
null set, hence has integral zero. If $\int f=C<\infty$, comparison
with $m\mathbf1_{\{f=\infty\}}$ gives
$\mu(\{f=\infty\})\leq C/m$ for every positive integer $m$, so this
set is null. Finally, decompose each of two almost-everywhere equal
functions into its restrictions to a measurable null exceptional set
and its complement. The null-supported parts have integral zero and
the other parts agree. Additivity proves equality of their integrals.
All exceptional sets used here are measurable; no completion is needed.
\Done

\begin{theorem}[Fatou's lemma]
\label{thm:fatou}
For nonnegative measurable $f_n$,
\[
\int\liminf_n f_n\,d\mu\leq\liminf_n\int f_n\,d\mu.
\]
\end{theorem}
\Proof
The measurable functions $g_n=\inf_{k\geq n}f_k$ increase to the lower
limit. For every $k\geq n$, monotonicity gives $\int g_n\leq\int f_k$,
hence $\int g_n\leq\inf_{k\geq n}\int f_k$. Take increasing limits
and apply monotone convergence to $g_n$.
\Done

For a real measurable $f$ with $\int|f|<\infty$, define
$\int f=\int f^+-\int f^-$, where $f^+=\max(f,0)$ and
$f^-=\max(-f,0)$. Both integrals are finite. This gives a linear
integral on such functions. First, $|f+g|\leq|f|+|g|$ and the
nonnegative additivity already proved show that $f+g$ is integrable.
Additivity then follows by integrating
\[
f^++g^++(f+g)^-=f^-+g^-+(f+g)^+
\]
and rearranging finite quantities; homogeneity follows by positive scaling
and interchange of positive and negative parts. Define the complex
integral by real and imaginary parts. In either case,
$|\int f|\leq\int|f|$: in the complex case choose a unit scalar making
$\int f$ real nonnegative and use $\operatorname{Re}(\lambda f)\leq|f|$.
The real case is the same estimate with a sign.

\begin{theorem}[Dominated convergence in norm]
\label{thm:dominated_convergence}
Let $f_n,f$ be measurable real or complex functions, with $f_n\to f$
almost everywhere and $|f_n|\leq g$ almost everywhere for some
nonnegative $g$ with $\int g<\infty$. Then $f$ and every $f_n$ are
integrable and
\[
\int|f_n-f|\,d\mu\longrightarrow0,
\qquad \int f_n\,d\mu\longrightarrow\int f\,d\mu.
\]
\end{theorem}
\Proof
Remove one measurable null set containing the convergence exception,
all the countably many domination exceptions, and $\{g=\infty\}$.
Set all functions to zero there; the resulting functions remain measurable.
The preceding null-invariance lemma applies first to their nonnegative
absolute values. For the modified functions $|f|\leq g$ everywhere,
so they are integrable; equality of the absolute-value integrals then
proves integrability of the original functions as well. Positive and
negative parts (or real and imaginary parts) show that their integrals
are unchanged.
Put $h_n=|f_n-f|$. Then $0\leq h_n\leq2g$ and $h_n\to0$.
Fatou applied to $2g-h_n$ gives
\[
2\int g\leq\liminf_n\left(2\int g-\int h_n\right)
=2\int g-\limsup_n\int h_n.
\]
All quantities subtracted are finite, so $\limsup_n\int h_n\leq0$.
Nonnegativity proves convergence to zero. The integral inequality
$|\int f_n-\int f|\leq\int h_n$ proves the final assertion.
\Done

\subsection{Functions and their almost-everywhere classes}
With this integral, the real space
$L^1(\mu)$ consists of equivalence classes of measurable functions
$f:\Omega\to\mathbb{R}$ satisfying
\[
\int_\Omega |f|\,d\mu<\infty,
\]
where two functions represent the same element when they agree almost
everywhere. Its norm is
\[
\|[f]\|_1=\int_\Omega |f|\,d\mu.
\]
The integral is unchanged by changing the representative on a null set.
Point evaluation $[f]\mapsto f(\omega)$, however, generally is not
well-defined on the quotient. This is a different issue from the book's
choice of a default value for an unevaluated class expression.

The complex version of $L^1$ uses complex measurable representatives
and the same absolute-value norm; all the integral estimates remain valid.
The scalar field is fixed when discussing either version as a vector space.

The space $L^\infty(\mu)$ consists of essentially bounded measurable
functions over that field, modulo the same equality almost everywhere.
Its norm is
\[
\|[g]\|_\infty=\inf\{C\geq0:|g|\leq C\text{ almost everywhere}\}.
\]
The essential bound at this infimum is attained. If its value is $M$,
then $|g|\leq M+1/n$ almost everywhere for each positive integer $n$:
by the definition of infimum, some admissible bound is less than
$M+1/n$. Remove the union of these countably many measurable null
exceptional sets and let $n\to\infty$ to get $|g|\leq M$ there.
This also proves that norm zero means the zero almost-everywhere class.
Homogeneity and the triangle inequality follow by combining attained
bounds off the union of the exceptional sets. Thus the expression is
indeed a norm, independent of the representative; no completeness of
the measure is required. For
$g\in L^\infty(\mu)$ and $f\in L^1(\mu)$, the pairing
\[
\langle g,f\rangle=\int_\Omega gf\,d\mu
\]
is independent of representatives and satisfies
$|\langle g,f\rangle|\leqslant\|g\|_\infty\|f\|_1$.
It is this pairing, rather than pointwise evaluation, that enters the
weak-convergence statements below. The following arguments derive
completeness from the convergence theorems just proved.

\subsection{Convergence and completeness}
In this subsection we use Theorems~\ref{thm:monotone_convergence}
and~\ref{thm:fatou}. These integral properties do not alter the earlier
choices of class evaluation.
A Banach space is a normed vector space complete for its norm metric.
A Hilbert space is an inner-product space complete for the norm induced
by its inner product.
The completeness arguments below use countable choice to select
measurable representatives of a sequence of almost-everywhere equivalence
classes. Subsequence indices can be chosen as the least suitable natural
numbers, so their recursive selection needs no additional dependent-choice
principle. These assumptions concern the quotient-space proofs; the
preceding monotone and dominated convergence arguments concern already
given measurable functions.

\begin{lemma}[Tail integrals and small sets]
For $f\in L^1(\mu)$ on any positive measure space,
\[
\int_{\{|f|>M\}}|f|\,d\mu\longrightarrow0
\quad(M\to\infty).
\]
Consequently, for every $\varepsilon>0$ there is $\delta>0$ such that
$\mu(A)<\delta$ implies $\int_A|f|\,d\mu<\varepsilon$.
\end{lemma}
\Proof
The functions $|f|\mathbf1_{\{|f|\leqslant m\}}$ increase pointwise
to $|f|$ as the positive integers $m$ tend to infinity. Monotone
convergence and finiteness of $\int|f|$ imply that their complementary
tail integrals tend to zero. Monotonicity in the real parameter $M$
gives the stated limit. Choose $M>0$ with tail integral less than
$\varepsilon/2$, and put $\delta=\varepsilon/(2M)$. Splitting $A$
into the regions $|f|\leqslant M$ and $|f|>M$ gives
\[
\int_A|f|\,d\mu\leqslant M\mu(A)
 +\int_{\{|f|>M\}}|f|\,d\mu<\varepsilon.
\]
No finiteness assumption on $\mu(\Omega)$ enters this argument.
\Done

For $a>0$, the same integral comparison gives
$a\mu(\{|f|>a\})\leqslant\|f\|_1$. This inequality will turn
norm convergence into convergence in measure. Here convergence in
measure means $\mu(\{|f_n-f|>a\})\to0$ for every $a>0$; on an
infinite-measure space this is a global condition, not merely convergence
on each finite-measure subset.

\begin{lemma}[A subsequence converging almost everywhere]
\label{lem:measure_subsequence}
If measurable real functions $f_n$ converge in measure to a measurable
real function $f$, then some subsequence converges to $f$ almost everywhere.
The measure space need not have finite measure.
\end{lemma}
\Proof
Choose increasing indices $n_k$ with
$\mu(\{|f_{n_k}-f|>2^{-k}\})<2^{-k}$. If $E_k$ denotes this
exceptional set, countable subadditivity gives
\[
\mu\left(\bigcap_{m\geqslant1}\bigcup_{k\geqslant m}E_k\right)
\leqslant\sum_{k\geqslant m}2^{-k}
\quad\text{for every }m.
\]
Thus the limsup of the $E_k$ is a measurable null set. Outside it,
$|f_{n_k}-f|\leqslant2^{-k}$ eventually, which proves the assertion.
\Done

\begin{theorem}[Completeness of $L^1$]
For any positive measure space, $L^1(\mu)$ is a Banach space.
\end{theorem}
\Proof
Let $(f_n)$ be Cauchy in norm. Choose a subsequence $(f_{n_k})$ with
$\|f_{n_{k+1}}-f_{n_k}\|_1\leqslant2^{-k}$, and choose measurable
representatives. Monotone convergence applied to the partial sums gives
\[
\int_\Omega\sum_{k=1}^\infty
 |f_{n_{k+1}}-f_{n_k}|\,d\mu\leqslant1.
\]
Consequently this series is finite outside a measurable null set.
There the telescoping series defines a finite limit $f$; put $f=0$
on the exceptional set. This gives a measurable function, even when
the measure is not complete, since the exceptional set is measurable.
The pointwise bound
\[
|f|\leqslant |f_{n_1}|+
\sum_{k=1}^\infty|f_{n_{k+1}}-f_{n_k}|
\quad\text{almost everywhere}
\]
shows that $f\in L^1$. Moreover,
\[
\|f-f_{n_j}\|_1\leqslant\sum_{k=j}^\infty2^{-k}\longrightarrow0.
\]
The Cauchy property and the triangle inequality now imply that the
whole sequence converges to $f$ in norm.
\Done

\begin{theorem}[Completeness of $L^\infty$]
For any positive measure space, $L^\infty(\mu)$ is a Banach space.
\end{theorem}
\Proof
Given a norm-Cauchy sequence, choose a subsequence with successive
norm distances at most $2^{-k}$. For measurable representatives,
the bounds $|f_{n_{k+1}}-f_{n_k}|\leqslant2^{-k}$ and the essential
bound for $f_{n_1}$ hold simultaneously off a measurable null set:
only countably many exceptional sets have been removed. The telescoping
series converges uniformly on that complement to a bounded function.
Extending the limit by zero on the exceptional set gives a measurable,
essentially bounded function $f$ with
$\|f-f_{n_j}\|_\infty\leqslant\sum_{k\geqslant j}2^{-k}$.
As before, the original Cauchy sequence has the same norm limit.
\Done

\subsection{The Hilbert space used by truncation}
For the real or complex scalar field, define $L^2(\mu)$ to be the
almost-everywhere classes of measurable functions with $\int|f|^2<\infty$.
Put
\[
\langle f,g\rangle=\int f\overline g\,d\mu,
\qquad \|f\|_2=\left(\int|f|^2\,d\mu\right)^{1/2}.
\]
Conjugation is omitted over the reals. The same measurable-null-set
convention applies as for $L^1$.
Here the angle brackets denote the $L^2$ inner product; over complex
scalars it conjugates the second factor. The earlier $L^\infty$--$L^1$
pairing $\int gf\,d\mu$ is bilinear and uses no conjugation.

\begin{lemma}[Cauchy--Schwarz and the $L^2$ norm]
The pairing is a well-defined inner product, linear in its first
argument, and
\[
|\langle f,g\rangle|\leq\|f\|_2\|g\|_2,
\qquad \|f+g\|_2\leq\|f\|_2+\|g\|_2.
\]
\end{lemma}
\Proof
The bound $2|f\overline g|\leq|f|^2+|g|^2$ proves integrability of
the product. Linearity, conjugate symmetry, and independence of
representatives follow from the integral. A nonnegative function has
zero integral precisely when it is zero almost everywhere, so the
pairing is positive definite on the quotient. If $\|g\|_2>0$, expand
$0\leq\|f-cg\|_2^2$ with
$c=\langle f,g\rangle/\|g\|_2^2$ to obtain Cauchy--Schwarz. If
$\|g\|_2=0$, its equivalence class is zero and the result is immediate.
Expanding $\|f+g\|_2^2$ and applying Cauchy--Schwarz gives the triangle
inequality. Closure under sums follows already from
$|f+g|^2\leq2|f|^2+2|g|^2$.
\Done

\begin{theorem}[Completeness of $L^2$]
\label{thm:L2_complete}
For any positive measure space, $L^2(\mu)$ is a Hilbert space.
\end{theorem}
\Proof
For a Cauchy sequence choose a subsequence with differences
$d_k=f_{n_{k+1}}-f_{n_k}$ satisfying $\|d_k\|_2\leq2^{-k}$,
and choose measurable representatives. The triangle inequality gives
\[
\left\|\sum_{k=j}^N|d_k|\right\|_2
\leq\sum_{k=j}^N2^{-k}.
\]
By monotone convergence applied to the squares, the measurable
nonnegative function $H_j=\sum_{k=j}^\infty|d_k|$ satisfies
\[
\int H_j^2\,d\mu\leq\left(\sum_{k=j}^\infty2^{-k}\right)^2.
\]
In particular $H_1$ is finite outside a measurable null set. The
corresponding telescoping series converges there to a finite measurable
function $f$; define $f=0$ on the exceptional set.
The bound $|f|\leq|f_{n_1}|+H_1$ almost everywhere proves $f\in L^2$.
Also $|f-f_{n_j}|\leq H_j$ almost everywhere, hence
$\|f-f_{n_j}\|_2\leq\sum_{k\geq j}2^{-k}\to0$.
The Cauchy property supplies convergence of the whole sequence. No
finiteness or completeness assumption on the measure is needed.
\Done

\ifdefined\qwertyflag
\undef\qwertyflag
\end{document}
\fi
